\section{Concrete historical analysis of the cycle}

\begin{frame}
  \frametitle{Cycle dating and phasing}
  Identifying the periods and phases of the industrial cycle, the crises of the capitalist economy are necessary to distinguish:

  \vspace{10pt}

  \begin{columns}[T,onlytextwidth]
    \begin{column}{0.48\textwidth}
      \textbf{In scale:}
      \begin{itemize}
        \item \textit{general overproduction} – whole national economy
        \item \textit{partial overproduction} – individual industries or sectors of the economy
      \end{itemize}
    \end{column}

    \begin{column}{0.48\textwidth}
      \textbf{In nature:}
      \begin{itemize}
        \item \textit{cyclical} – caused by immanent contradictions and reproduction of fixed capital
        \item \textit{intermediate} – other and mostly exogenous factors
      \end{itemize}
    \end{column}
  \end{columns}

  \begin{block}{Proposition}
    Industrial cycle dating and phasing is based on cyclical crises of general overproduction.
  \end{block}

  \begin{alertblock}{Important}
    Only quantitative data with qualitative analysis of historical conditions comparison makes distinguishing the cyclical crisis of general overproduction from non-cyclical recessions, sectoral imbalances and externally caused economic shocks possible.
  \end{alertblock}
\end{frame}

\begin{frame}[t]
  \frametitle{Cycle periods and phases (1)\footnotetext{Source: FRED, NBER}}
  \begin{center}
    \input{Plots/indpro-phases-general.tex} % Период и фазы цикла
  \end{center}
\end{frame}

\begin{frame}
  \frametitle{Cycle periods and phases (2)}
  \begin{alertblock}{Finding}
    The standard sequence \textquotedbl{}crisis $\to$ depression $\to$ recovery $\to$ growth\textquotedbl{}, reflecting the internal logic of the movement of the industrial cycle, may undergo some changes in the real dynamics of the capitalist economy.
  \end{alertblock}

  There are two forms of these deviations:

  \begin{columns}[T,onlytextwidth]
    \begin{column}{0.48\textwidth}
      \begin{itemize}
        \item \textit{intermediate crisis} – occurs under the influence of exogenous factors during the growth phase, after which it can either continue or end\tablefootnotetext{For example, during the crises of 1973 and 1945 respectively}
      \end{itemize}
    \end{column}

    \begin{column}{0.48\textwidth}
      \begin{itemize}
        \item \textit{intermediate recovery} – occurs during the crisis phase under the influence of a short-term effect of stimulating government policy, after which the crisis phase continues\tablefootnotetext{For example, during the crisis periods of 1957-1961 and 1980-1982}
      \end{itemize}
    \end{column}
  \end{columns}
\end{frame}

\begin{frame}[t]
  \frametitle{Cycle periods and phases duration (1)\footnotetext{Source: FRED, NBER}}
  \begin{center}
    % Данные
	\pgfplotstableread[col sep=tab]{Data/phases.tsv} \phases % Продолжительность фаз цикла и промежутков

% Область построения
\begin{tikzpicture}
    \begin{axis}[
        plotstyle-1,
        height=0.79\textheight,
        width=1.05\textwidth,
        ybar stacked,
        bar width=10pt,
        xtick=data,
        x tick label style={
            anchor=east,
        },
        xticklabels={
            1920-1929,
            1929-1937,
            1937-1948,
            1948-1957,
            1957-1969,
            1969-1980,
            1980-1990,
            1990-2001,
            2001-2007,
            2007-2020,
            2020-20??
        },
        %scale only axis,
        %height=0.65\textheight,
        %width=1.30\textwidth,
        ytick distance=10,
        xmin=0.5, xmax=11.5,
        ymin=0, ymax=150,
        ylabel={months},
        legend columns=-1,
        tick style={draw=none}
    ]

% Средняя продолжительность периода
\draw[dashed, black, line width=0.5pt] (axis cs:0.5,120.1) -- (axis cs:11.5,120.1);

% Средняя продолжительность фаз
    % Кризис
    \filldraw[
        fill={rgb,255:red,128;green,0;blue,0},
        draw=black,
    ]
    (axis cs:0.5,0) rectangle (axis cs:0.55,13.4);

    % Депрессия
    \filldraw[
        fill={rgb,255:red,61;green,23;blue,199},
        draw=black,
    ]
    (axis cs:0.5,13.4) rectangle (axis cs:0.55,16.7);

    % Оживление
    \filldraw[
        fill={rgb,255:red,51;green,153;blue,153},
        draw=black,
    ]
    (axis cs:0.5,16.7) rectangle (axis cs:0.55,40);

    % Подъем
    \filldraw[
        fill={rgb,255:red,220;green,173;blue,118},
        draw=black,
    ]
    (axis cs:0.5,40) rectangle (axis cs:0.55,104.1);

    % Промежутки
    \filldraw[
        fill=gray, opacity=0.5,
        draw=black,
    ]
    (axis cs:0.5,104.1) rectangle (axis cs:0.55,120.1);

% Прогноз последней фазы подъема
\filldraw[
    fill={rgb,255:red,220;green,173;blue,118},
    draw=black,
]
(axis cs:10.83,63) rectangle (axis cs:11.17,74) node[midway] {11};

\draw[draw=black, line width=0.5pt] % Интервал погрешности
    (axis cs:11,96) -- (axis cs:11,144);

\draw[draw=black, line width=0.5pt] % Нижняя засечка
    (axis cs:10.94,96) -- (axis cs:11.06,96);

\draw[draw=black, line width=0.5pt] % Верхняя засечка
    (axis cs:10.94,144) -- (axis cs:11.06,144);

\fill (axis cs:11,120.1) circle[radius=1pt]; % Круг и оценка

\filldraw[
    fill=white,
    draw=black,
]
(axis cs:10.83,74) rectangle (axis cs:11.17,120.1) node[midway, yshift=-13pt] {$\widetilde{46}$};

% Кризис
\addplot[fill={rgb,255:red,128;green,0;blue,0},
    forget plot,
    nodes near coords={
        \pgfmathfloattofixed{\pgfplotspointmeta}%
        \let\labelvalue\pgfmathresult%
        \pgfmathtruncatemacro{\showlabel}{\labelvalue >= 5}%
        \ifnum\showlabel=1%
            \pgfmathprintnumber{\labelvalue}%
        \fi%
    },
    nodes near coords style={/pgf/number format/.cd,precision=0},
] table[
    x=№,
    y=Кризис,
] {\phases};

% Промежутки К
\addplot[fill={rgb,255:red,128;green,0;blue,0}, opacity=0.6,
    forget plot,
    nodes near coords,
    nodes near coords style={/pgf/number format/.cd,precision=0},
] table[
    x=№,
    y=Промежутки К,
] {\phases};

% Депрессия
\addplot[fill={rgb,255:red,61;green,23;blue,199},
    forget plot,
    nodes near coords={
        \pgfmathfloattofixed{\pgfplotspointmeta}%
        \let\labelvalue\pgfmathresult%
        \pgfmathtruncatemacro{\showlabel}{\labelvalue >= 5}%
        \ifnum\showlabel=1%
            \pgfmathprintnumber{\labelvalue}%
        \fi%
    },
    nodes near coords style={/pgf/number format/.cd,precision=0},
] table[
    x=№,
    y=Депрессия,
    meta=Депрессия,
] {\phases};

% Оживление
\addplot[fill={rgb,255:red,51;green,153;blue,153},
    forget plot,
    nodes near coords,
    nodes near coords style={/pgf/number format/.cd,precision=0},
] table[
    x=№,
    y=Оживление,
] {\phases};

% Подъем
\addplot[fill={rgb,255:red,220;green,173;blue,118},
    forget plot,
    nodes near coords,
    nodes near coords style={/pgf/number format/.cd,precision=0},
] table[
    x=№,
    y=Подъем,
] {\phases};

% Промежутки П
\addplot[fill={rgb,255:red,220;green,173;blue,118}, opacity=0.6,
    forget plot,
    nodes near coords,
    nodes near coords style={/pgf/number format/.cd,precision=0},
] table[
    x=№,
    y=Промежутки П,
] {\phases};

% Легенда
\addlegendimage{
    legend image code/.code={
        \fill[fill={rgb,255:red,128;green,0;blue,0}, opacity=0.8]
            (0cm,-0.1cm) rectangle (0.16cm,0.1cm);
        \fill[fill={rgb,255:red,61;green,23;blue,199}, opacity=0.8]
            (0.22cm,-0.1cm) rectangle (0.38cm,0.1cm);
        \fill[fill={rgb,255:red,51;green,153;blue,153}, opacity=0.8]
            (0.44cm,-0.1cm) rectangle (0.60cm,0.1cm);
        \fill[fill={rgb,255:red,220;green,173;blue,118}, opacity=0.8]
            (0.66cm,-0.1cm) rectangle (0.82cm,0.1cm);
    }
}
\addlegendentry{Phases CDRG}

\addlegendimage{
        legend image code/.code={
            \filldraw[fill={rgb,255:red,128;green,0;blue,0}, draw=black]
                (0.24cm,-0.30cm) rectangle (0.36cm,-0.233cm);
            \filldraw[fill={rgb,255:red,61;green,23;blue,199}, draw=black]
                (0.24cm,-0.233cm) rectangle (0.36cm,-0.217cm);
            \filldraw[fill={rgb,255:red,51;green,153;blue,153}, draw=black]
                (0.24cm,-0.217cm) rectangle (0.36cm,-0.100cm);
            \filldraw[fill={rgb,255:red,220;green,173;blue,118}, draw=black]
                (0.24cm,-0.100cm) rectangle (0.36cm,0.220cm);
            \filldraw[
                fill=gray,
                fill opacity=0.5,
                draw=black,
                draw opacity=1,
            ] (0.24cm,0.220cm) rectangle (0.36cm,0.300cm);
        }
    }
    \addlegendentry{Average phases}

\addlegendimage{
    legend image code/.code={
        \draw[dashed, black, line width=0.5pt]
            (0cm,0cm) -- (0.6cm,0cm);
    }
}
\addlegendentry{Average period}

\addlegendimage{
    legend image code/.code={
        \fill[fill=white, opacity=0.8]
            (0cm,-0.1cm) rectangle (0.16cm,0.1cm);
    }
}
\addlegendentry{Forecast}

\addlegendimage{
    legend image code/.code={
        \draw[black, line width=0.5pt]
            (0cm,0cm) -- (0.6cm,0cm);
        \draw[black, line width=0.5pt]
            (0cm,-0.1cm) -- (0cm,0.1cm);
        \draw[black, line width=0.5pt]
            (0.6cm,-0.1cm) -- (0.6cm,0.1cm);
        \fill[black] (0.3cm,0cm) circle[radius=1pt];
    }
    }
    \addlegendentry{Error estimate}
            
        \end{axis}
\end{tikzpicture} % % Продолжительность периодов и фаз цикла
  \end{center}
\end{frame}

\begin{frame}
  \frametitle{Cycle periods and phases duration (2)}
  Cycle periods and phases duration within the study period was in months:

  \vspace{10pt}

  \rowcolors{2}{gray!10}{white}
  \begin{tabularx}{\textwidth}{|>{\raggedright\arraybackslash\hyphenpenalty=10000\exhyphenpenalty=10000}X|r|r|r|r|r|}
    \hline
    \cellcolor{rgb,255:red,15;green,35;blue,84}\textcolor{white}{Duration} & Period & \cellcolor{rgb,255:red,128;green,0;blue,0}\textcolor{white}{Crisis} & \cellcolor{rgb,255:red,61;green,23;blue,199}\textcolor{white}{Depression} & \cellcolor{rgb,255:red,51;green,153;blue,153}\textcolor{white}{Recovery} & \cellcolor{rgb,255:red,220;green,173;blue,118}\textcolor{white}{Growth} \\ \hline
    Average & 120.1 & 16.6 & 3.3 & 23.0 & 79.5 \\ \hline
    Adjusted & 120.1 & 13.4 & 3.3 & 23.0 & 64.1 \\ \hline
    Maximum & 148.0 & 35.0 & 8.0 & 60.0 & 107.0 \\ \hline
    Minimum & 81.0 & 2.0 & 1.0 & 5.0 & 30.0 \\ \hline
  \end{tabularx}

  \begin{block}{Observation}
    Duration of the depression phase of the industrial cycle decreased from 6-8 months to 1-2 months from 1919 to 2026, which can be attributed to the development of government incentive policies.
  \end{block}

  \begin{alertblock}{Finding}
    Empirically obtained data confirm the proposition about 10 years average duration of the industrial cycle period, corresponding the period of physical depreciation of fixed capital in the field of industrial production.
  \end{alertblock}
\end{frame}

\begin{frame}
  \frametitle{Current industrial cycle}
  The current industrial cycle in the US economy began after the crisis slump in production in 2020 has a number of features:
  \begin{itemize}
    \item An extremely long recovery phase after the 2020 crisis of 60 months with an average duration of 23 months;
    \item The economy is formally in growth phase relative to the last crisis, according to NBER dating, but in fact the pre-crisis maximum of 2018 of 104 points has not been overcome yet;
    \item The pre-crisis level of 2007 has hardly been overcome only twice in 2014 and 2018 respectively, and at the moment the index value is still below the maximum of 2007, indicating the stagnation of industrial production.
  \end{itemize}

  \begin{alertblock}{Hypothesis}
    The next crisis of general overproduction in the United States based on the frequency of the industrial cycle identified in this paper all other things being equal may unfold in the period from 2028 to 2032.

    \vspace{10pt}

    A potential reason for a new cyclical crisis may be the depreciation of capital in the artificial intelligence market due to an inflated financial bubble based on exagerated expectations of future profits.
  \end{alertblock}
\end{frame}